% title: CT18Z nominal fits at reco and gen level, with and without the lattice, cold trust-constr
% slug: ct18z-nominal-fits
% status: done
% period: 2026-10-15 – 2026-10-20
% updated: 2026-10-20
% covers: 2026-10-20
%
% BUILD: scripts/summary_pdf.py ct18z-nominal-fits  (on the host; pdflatex is not in the container)
\documentclass[10pt,a4paper]{article}

% ---- page and fonts ----
\usepackage[a4paper,top=16mm,bottom=20mm,left=18mm,right=18mm,footskip=9mm]{geometry}
\usepackage[T1]{fontenc}
\usepackage[utf8]{inputenc}
\usepackage{lmodern}
\usepackage{microtype}
\usepackage{amsmath,amssymb}
\usepackage{graphicx}
\usepackage{booktabs}
\usepackage[table]{xcolor}
\usepackage[font=small,labelfont={bf,sf},margin=6pt,skip=4pt]{caption}
\usepackage{subcaption}
\usepackage{enumitem}
\usepackage{titlesec}
\usepackage{fancyhdr}
\usepackage[section,below]{placeins}   % a result's figure stays inside its section
\usepackage{tikz}                      % for a schematic / pipeline figure (Section 2)
\usetikzlibrary{arrows.meta,positioning}
\usepackage{url}
% siunitx is not installed on the submit host's TeX Live; these fallbacks cover the usual
% \num{0.547}, \SI{2.3}{\GeV}, \SI{0.14}{\percent} (no e-notation: write 1.2\times10^{-3}).
\IfFileExists{siunitx.sty}{\usepackage{siunitx}}{}
\providecommand\num[1]{\ensuremath{#1}}
\providecommand\si[1]{\ensuremath{\mathrm{#1}}}
\providecommand\SI[2]{\ensuremath{#1\,\mathrm{#2}}}
\providecommand\GeV{GeV}\providecommand\MeV{MeV}\providecommand\TeV{TeV}
\providecommand\percent{\%}\providecommand\pb{pb}\providecommand\fb{fb}
\definecolor{sumblue}{HTML}{1B5FAE}
\definecolor{sumrule}{HTML}{2A78D6}
\definecolor{sumbox}{HTML}{EEF4FB}
\definecolor{sumgray}{HTML}{5A5A5A}
\definecolor{sumwarn}{HTML}{8A5A00}
\usepackage[colorlinks=true,linkcolor=sumblue,urlcolor=sumblue,citecolor=sumblue,
            pdfborder={0 0 0}]{hyperref}

% unicode copied from a logbook would otherwise stop pdflatex
\DeclareUnicodeCharacter{03B1}{\ensuremath{\alpha}}
\DeclareUnicodeCharacter{03B2}{\ensuremath{\beta}}
\DeclareUnicodeCharacter{03B3}{\ensuremath{\gamma}}
\DeclareUnicodeCharacter{03B4}{\ensuremath{\delta}}
\DeclareUnicodeCharacter{03B5}{\ensuremath{\epsilon}}
\DeclareUnicodeCharacter{03B7}{\ensuremath{\eta}}
\DeclareUnicodeCharacter{03B8}{\ensuremath{\theta}}
\DeclareUnicodeCharacter{03BA}{\ensuremath{\kappa}}
\DeclareUnicodeCharacter{03BB}{\ensuremath{\lambda}}
\DeclareUnicodeCharacter{03BC}{\ensuremath{\mu}}
\DeclareUnicodeCharacter{03BD}{\ensuremath{\nu}}
\DeclareUnicodeCharacter{03C0}{\ensuremath{\pi}}
\DeclareUnicodeCharacter{03C1}{\ensuremath{\rho}}
\DeclareUnicodeCharacter{03C3}{\ensuremath{\sigma}}
\DeclareUnicodeCharacter{03C4}{\ensuremath{\tau}}
\DeclareUnicodeCharacter{03C6}{\ensuremath{\phi}}
\DeclareUnicodeCharacter{03C7}{\ensuremath{\chi}}
\DeclareUnicodeCharacter{03C8}{\ensuremath{\psi}}
\DeclareUnicodeCharacter{03C9}{\ensuremath{\omega}}
\DeclareUnicodeCharacter{0393}{\ensuremath{\Gamma}}
\DeclareUnicodeCharacter{0394}{\ensuremath{\Delta}}
\DeclareUnicodeCharacter{039B}{\ensuremath{\Lambda}}
\DeclareUnicodeCharacter{03A3}{\ensuremath{\Sigma}}
\DeclareUnicodeCharacter{03A6}{\ensuremath{\Phi}}
\DeclareUnicodeCharacter{03A9}{\ensuremath{\Omega}}
\DeclareUnicodeCharacter{2212}{\ensuremath{-}}
\DeclareUnicodeCharacter{2248}{\ensuremath{\approx}}
\DeclareUnicodeCharacter{2264}{\ensuremath{\leq}}
\DeclareUnicodeCharacter{2265}{\ensuremath{\geq}}
\DeclareUnicodeCharacter{2260}{\ensuremath{\neq}}
\DeclareUnicodeCharacter{226A}{\ensuremath{\ll}}
\DeclareUnicodeCharacter{226B}{\ensuremath{\gg}}
\DeclareUnicodeCharacter{223C}{\ensuremath{\sim}}
\DeclareUnicodeCharacter{221D}{\ensuremath{\propto}}
\DeclareUnicodeCharacter{221E}{\ensuremath{\infty}}
\DeclareUnicodeCharacter{221A}{\ensuremath{\surd}}
\DeclareUnicodeCharacter{2202}{\ensuremath{\partial}}
\DeclareUnicodeCharacter{2192}{\ensuremath{\rightarrow}}
\DeclareUnicodeCharacter{2190}{\ensuremath{\leftarrow}}
\DeclareUnicodeCharacter{21D2}{\ensuremath{\Rightarrow}}
\DeclareUnicodeCharacter{2194}{\ensuremath{\leftrightarrow}}
\DeclareUnicodeCharacter{2208}{\ensuremath{\in}}
\DeclareUnicodeCharacter{2026}{\ldots}
\DeclareUnicodeCharacter{2713}{\ensuremath{\checkmark}}
\DeclareUnicodeCharacter{2080}{\ensuremath{_0}}
\DeclareUnicodeCharacter{2081}{\ensuremath{_1}}
\DeclareUnicodeCharacter{2082}{\ensuremath{_2}}
\DeclareUnicodeCharacter{2083}{\ensuremath{_3}}
\DeclareUnicodeCharacter{2084}{\ensuremath{_4}}
\DeclareUnicodeCharacter{2085}{\ensuremath{_5}}
\DeclareUnicodeCharacter{2086}{\ensuremath{_6}}
\DeclareUnicodeCharacter{2087}{\ensuremath{_7}}
\DeclareUnicodeCharacter{2088}{\ensuremath{_8}}
\DeclareUnicodeCharacter{2089}{\ensuremath{_9}}
\DeclareUnicodeCharacter{2070}{\ensuremath{^{0}}}
\DeclareUnicodeCharacter{2074}{\ensuremath{^{4}}}
\DeclareUnicodeCharacter{2075}{\ensuremath{^{5}}}
\DeclareUnicodeCharacter{2076}{\ensuremath{^{6}}}
\DeclareUnicodeCharacter{2077}{\ensuremath{^{7}}}
\DeclareUnicodeCharacter{2078}{\ensuremath{^{8}}}
\DeclareUnicodeCharacter{2079}{\ensuremath{^{9}}}
\DeclareUnicodeCharacter{207B}{\ensuremath{^{-}}}
\DeclareUnicodeCharacter{207A}{\ensuremath{^{+}}}

% ---- layout ----
\setlength{\parindent}{0pt}
\setlength{\parskip}{0.55ex plus 0.2ex}
\setlist{itemsep=1pt,topsep=2pt,parsep=0pt,leftmargin=1.4em}
\titleformat{\section}{\sffamily\large\bfseries}{\thesection}{0.6em}{}
\titleformat{\subsection}{\sffamily\normalsize\bfseries}{\thesubsection}{0.6em}{}
\titlespacing*{\section}{0pt}{2.0ex plus 0.6ex minus 0.3ex}{0.7ex}
\titlespacing*{\subsection}{0pt}{1.4ex plus 0.4ex minus 0.2ex}{0.4ex}
\renewcommand{\topfraction}{0.9}
\renewcommand{\bottomfraction}{0.8}
\renewcommand{\textfraction}{0.07}
\renewcommand{\floatpagefraction}{0.75}
\setlength{\textfloatsep}{10pt plus 3pt minus 3pt}
\setlength{\floatsep}{8pt plus 2pt minus 2pt}
\setlength{\intextsep}{8pt plus 2pt minus 2pt}

% ---- metadata: passed in by scripts/summary_pdf.py from the "% key:" lines ----
\makeatletter
\providecommand\SumTitle{\textcolor{red}{(build with scripts/summary\_pdf.py)}}
\providecommand\SumSlug{?}
\providecommand\SumStudy{?}       % the study slug (the parent study, for a task summary)
\providecommand\SumViewPath{?}    % <slug> or <study>/<task>: the viewer anchor
\providecommand\SumStatus{?}
\providecommand\SumPeriod{?}
\providecommand\SumUpdated{?}
\providecommand\SumCovers{?}
\begingroup\catcode`\#=12 \catcode`\~=12
\gdef\Sum@web{https://submit.mit.edu/~lavezzo/alphaS/studies/#}
\endgroup

% ---- macros for the body ----
% \taskref{261008-task}: cite the task a number comes from; links to its viewer page.
\newcommand\taskref[1]{{\small\sffamily[\href{\Sum@web\SumStudy/#1}{#1}]}}
% \logref{2026-07-02}: cite a dated entry of a single-layer study (no task dirs).
\newcommand\logref[1]{{\small\sffamily[\href{\Sum@web\SumViewPath:logbook}{log #1}]}}
% \viewerlink{<study>/<task>}{text}: any other page of the viewer.
\newcommand\viewerlink[2]{\href{\Sum@web#1}{#2}}
\newcommand\logbooklink[1]{\href{\Sum@web\SumViewPath:logbook}{#1}}
% \code{path/or_flag}: typewriter, breaks at / and _, underscores need no escaping.
\newcommand\code{\nolinkurl}
% \caveat{...}: a comparability caveat set off before the number it qualifies.
\newcommand\caveat[1]{{\color{sumwarn}\small\sffamily\textbf{Caveat.}}~{\small\itshape #1}}
% physics shorthands
\newcommand\alphas{\ensuremath{\alpha_\mathrm{s}}}
\newcommand\dalphas{\ensuremath{\Delta\alpha_\mathrm{s}}}
\newcommand\pT{\ensuremath{p_\mathrm{T}}}
\newcommand\ptll{\ensuremath{p_\mathrm{T}^{\ell\ell}}}
\newcommand\qT{\ensuremath{q_\mathrm{T}}}
\newcommand\yll{\ensuremath{y^{\ell\ell}}}

% the answer box: \begin{answer} ... \end{answer}, one or two sentences with the key number
\newsavebox\Sum@box
\newenvironment{answer}
  {\par\vspace{1.2ex}\noindent\begin{lrbox}{\Sum@box}%
   \begin{minipage}{\dimexpr\linewidth-3pt-2\fboxsep\relax}\vspace{1pt}%
   {\sffamily\small\bfseries\color{sumblue}Answer}\par\vspace{1pt}}
  {\vspace{1pt}\end{minipage}\end{lrbox}%
   \hbox{{\color{sumrule}\vrule width 3pt}\colorbox{sumbox}{\usebox{\Sum@box}}}%
   \par\vspace{0.8ex}}

\newcommand\SummaryTitleBlock{%
  {\sffamily\footnotesize\color{sumgray}%
   \alphas\ analysis\,\textperiodcentered\,study summary\hfill status: \SumStatus\par}%
  \vspace{3pt}%
  {\sffamily\LARGE\bfseries\SumTitle\par}%
  \vspace{5pt}%
  {\small\color{sumgray}%
   \texttt{\SumViewPath}\enspace\textperiodcentered\enspace period \SumPeriod
   \enspace\textperiodcentered\enspace written \SumUpdated
   \enspace\textperiodcentered\enspace covers logbook through \textbf{\SumCovers}\par
   Full record: \logbooklink{logbook}
   \enspace\textperiodcentered\enspace web page: \href{\Sum@web\SumViewPath}{\nolinkurl{submit.mit.edu/~lavezzo/alphaS/studies}}%
   \enspace\textperiodcentered\enspace plots in each task directory\par}%
  \vspace{4pt}{\color{sumgray}\hrule height 0.4pt}\vspace{2pt}}

\pagestyle{fancy}
\fancyhf{}
\renewcommand{\headrulewidth}{0pt}
\fancyfoot[L]{\sffamily\scriptsize\color{sumgray}\alphas\ study summary\enspace\textperiodcentered\enspace\texttt{\SumViewPath}\enspace\textperiodcentered\enspace covers logbook through \SumCovers}
\fancyfoot[R]{\sffamily\scriptsize\color{sumgray}\thepage}
\makeatother

\hypersetup{pdftitle={\SumTitle}}

\begin{document}
\thispagestyle{fancy}
\SummaryTitleBlock
\begin{answer}
Four cold trust-constr fits to real CT18Z data with the physical NP region as hard constraints (reco $\ptll\times\yll$
and unfolded gen-level $\sigma_\text{UL}$, each without and with the ASWZ lattice CS-kernel $\chi^2$) all converge, and each
level lands on the same NP faces. The lattice costs $2\Delta\mathrm{NLL}=13.5$ (reco) / 11.0 (gen) for two CS
parameters, moves the NP damping from the CS kernel to the TMD factor, shifts $\alphas$ by $+0.21\,\sigma$ / $+0.35\,\sigma$
and widens $\sigma(\alphas)$ by $+5\,\%$ / $+1\,\%$ (to $1.14$ / $1.18\times10^{-3}$); gen$-$reco is $-0.61\,\sigma$ (Z only) and $-0.42\,\sigma$
(with lattice). Caveats: $\alphas$ blinded (differences only); the full fits are acceptable but the \ptll\ (\ptll\
projection p = 0.02--0.25\,\%) and gen $p_\mathrm{T}^V$ shapes are not, a characterised mis-modelling left without a remedy;
the lattice fits sit on a weakly binding face, so their $\alphas$ interval needs a likelihood scan.
\end{answer}

% ============================================================================================
\section{Motivation and question}

The study went back to the nominal physics fits: CT18Z at the nominal order, without (``NOLAT'') and with
(``LAT'') the lattice information on the Collins--Soper (CS) kernel, now \emph{cold} (no warm start) and with the
physical NP region as hard constraints instead of a penalty wall. The questions (2026-10-15) were: (1) does the minimisation work cold; (2) what does the lattice do to the fit and the NP parameters;
(3) is the fit good. On 2026-10-17 three were added: (4) are these the same minima as the older trust-krylov + wall fits
on which earlier conclusions rest; (5) the same pair at gen level on the unfolded data, the gen-level half of the same four-way comparison
\taskref{261017-gen-cold-fits}; (6) why the \ptll\
projection of the GoF is poor.

% ============================================================================================
\section{Setup and method}

\textbf{Reco fits.} Real data, card \code{260916_Z_2D_card_adcorr/ZMassDilepton_ptll_yll_adexclpdf} (2D
$\ptll\times\yll$, 780 bins, Poisson, BB-lite), the N$^3$LL+NNLO SCETlib AD cache \code{pdf62_y35_260921_y25} (the
770-bin $|Y|\le2.5$ subset of \code{merged_full_bin0xzero}), \code{SCETlibADParamModel}; 3720 parameters
\taskref{261015-cold-tc-fits}. \textbf{Gen fits.} The $\sigma_\text{UL}$ of the iteration-0 real-data unfolding (20
$p_\mathrm{T}^V\times10$ $|Y^V|$ = 200 bins), covariance $\chi^2$, 6 card nuisances; the same cache folded onto the gen bins
(\code{gen_level=1}), same model defaults, constraints and lattice settings \taskref{261017-gen-cold-fits}. Unfolding-level systematics enter the
gen fits only through the data covariance.

\textbf{Common to all four.} No NP $\lambda$ carries a prior (the CS $\lambda_{2\nu},\lambda_{4\nu}$ were made prior-free for this
study), so LAT$-$NOLAT is exactly the lattice term (\code{LatticeCSTerm}, ASWZ points, frozen perturbative kernel,
$\tau=8$). \code{NPDomainConstraints} (margin 0) imposes six hard rows: $\lambda_{2\nu}\ge0$, $\lambda_{4\nu}\ge0$, the TMD
small-$b$ turn-on $L_2(Y)=\lambda_2+\delta\lambda_2Y^2\ge0$ and the TMD large-$b$ ``cubic'' $3\lambda_\infty^2\lambda_4+L_2^3\ge0$,
each at $|Y|=0$ and 2.5. Minimiser: scipy \code{trust-constr}, cold from the anchor, unpreconditioned (preconditioning
a cold start was found to hurt). Reco via \code{workflows/fitterAD.sh} with \code{--constraints} (and \code{--lattice}); gen through task
scripts with the same flags, since \code{fitterAD.sh} cannot run a gen card. All four fits use the same real-data
blinding frame (at gen forced with \code{--blindingFrame} \code{data}), so stored $\alphas$ differences across levels are physical.
The post-fit pass computes the face-held Hessian, traditional and global impacts, and the full and projected saturated
tests (\ptll\ at reco; $p_\mathrm{T}^V$ and $|Y^V|$ at gen).

% ============================================================================================
\section{Results}

\subsection{The four fits: same faces per configuration, a consistent lattice effect, gen below reco}

\caveat{Real data, $\alphas$ blinded: differences only, in units of the reference fit's $\sigma$ or in $10^{-3}$.
$\sigma(\alphas)$ is in $\alphas$ units $\times10^{-3}$ (twice rabbit's \code{alphaS} nuisance units), traditional total
with active faces held. Reco and gen are different likelihoods on the same events, not nested; NLLs compare only within a
level. Gen ``Data stat.'' is the whole unfolded-data
covariance (statistics, experimental, luminosity, response theory) and is not comparable to reco ``Data stat.''. LAT's
NLL and full saturated $q$ include the lattice $\chi^2$; ndof counts bins only.}

\begin{table}[htb]
  \centering\small
  \caption{The four nominal fits. Physical NP $\lambda$; ``face'' = parameter on a hard constraint, with its release
  $\Delta$NLL (the NLL gain if the constraint were removed). Impacts are traditional grouped, face held, in $10^{-3}$.
  $^*$Gen ``Data stat.'' = full unfolded covariance. \taskref{261015-cold-tc-fits} \taskref{261017-gen-cold-fits}}
  \label{tab:fits}
  \footnotesize\setlength{\tabcolsep}{3pt}
  \begin{tabular}{lllll}
    \toprule
     & reco NOLAT & reco LAT & gen NOLAT & gen LAT \\
    \midrule
    cold fit: iter., time & 62, 7.4\,h & 113, 4.8\,h & 57, 5.5\,h & 46, 1.6\,h \\
    lattice cost $2\Delta$NLL & & 13.51 & & 10.99 \\
    faces (release $\Delta$NLL) & $\lambda_{2\nu}{\ge}0$ (15.6); cubic$_{2.5}$ & $L_2(2.5)$ (0.46) & $\lambda_{2\nu}{\ge}0$ (14.1); cubic$_{2.5}$ (4.9) & $L_2(2.5)$ (0.34) \\
    full saturated $q$/ndf, $p$ & 739.4/774, 81\,\% & 753.0/774, 70\,\% & 226.4/194, 5.6\,\% & 237.4/194, 1.8\,\% \\
    proj.\ \ptll\ / $p_\mathrm{T}^V$ & 68.4/39, 0.25\,\% & 79.1/39, 0.02\,\% & 52.5/20, $10^{-4}$ & 60.3/20, $6\times10^{-6}$ \\
    proj.\ $|Y^V|$ & & & 16.6/10, 8.3\,\% & 16.4/10, 9.0\,\% \\
    $\sigma(\alphas)$ [$10^{-3}$] & 1.09 & 1.14 & 1.17 & 1.18 \\
    $\dalphas$ LAT$-$NOLAT [$\sigma_\text{NOLAT}$] & & $+0.21$ & & $+0.35$ \\
    $\dalphas$ gen$-$reco [$\sigma_\text{reco}$] & & & $-0.61$ & $-0.42$ \\
    CS $\lambda_{2\nu}$ / $\lambda_{4\nu}$ & 0 (face) / 0.045 & 0.040 / 0.011 & 0 (face) / 0.031 & 0.079 / 0.005 \\
    TMD $\lambda_2$ / $\lambda_4$ / $\delta\lambda_2$ & 0.115 / 0.000 / $-0.009$ & 0.046 / 0.072 / $-0.007$ & 0.171 / $-0.001$ / $-0.007$ & 0.018 / 0.042 / $-0.003$ \\
    impacts: NP (lead.\ sub-group) & 0.86 ($\gamma_\nu$ 0.83) & 0.93 ($F_\text{eff}$ 0.83) & 0.89 ($\gamma_\nu$ 0.85) & 0.91 ($F_\text{eff}$ 0.81) \\
    impacts: stat.\ / PDF / TNP & 0.79 / 0.71 / 0.65 & 0.86 / 0.79 / 0.77 & 1.05$^*$ / 0.77 / 0.69 & 1.09$^*$ / 0.79 / 0.77 \\
    \bottomrule
  \end{tabular}
\end{table}

\begin{figure}[htb]
  \centering
  \includegraphics[width=0.37\linewidth]{summary_figs/four_fits_alphas.pdf}
  \caption{Blinded $\alphas$ of the four fits relative to reco NOLAT, in $10^{-3}$, all in the same real-data blinding
  frame; the error bar is each fit's own face-held $\sigma(\alphas)$, \emph{not} the uncertainty of the difference (the
  fits share the events and are strongly correlated). Reco and gen are different likelihoods (card, histmaker,
  systematics treatment). Made from \code{compare.json} by \code{summary_figs/plot_four_fits.py}. \taskref{261017-gen-cold-fits}}
  \label{fig:alphas}
\end{figure}

\textbf{Convergence.} All four main fits reach trust-constr status 1 from the cold anchor (optimality $\le7\times10^{-9}$,
no violation) in 1.6--7.4\,h
\taskref{261015-cold-tc-fits} \taskref{261017-gen-cold-fits}. \textbf{Faces.} The configuration, not the level, decides the faces: without the
lattice the Z data put the CS small-$b$ coefficient on its bound ($\lambda_{2\nu}=0$, strongly binding) and the TMD
large-$b$ damping on its cubic face at $|Y|=2.5$; with the lattice the CS kernel moves onto the lattice side and only the
TMD small-$b$ face at $|Y|=2.5$ is active, weakly binding. \textbf{Lattice.} The lattice costs 13.5 (reco, $\approx3.2\sigma$
for 2 dof) and 11.0 (gen, $\approx2.9\sigma$), moves $\alphas$ up at both levels ($+0.21$ / $+0.35\,\sigma$) and makes
the NP sector's residual freedom move from the CS kernel ($\gamma_\nu$ group) to the TMD factor ($F_\text{eff}$ group); the NP
group stays the largest single source of $\sigma(\alphas)$ at reco. \textbf{Gen vs reco.} Gen lies below reco by
$0.66$ / $0.48\times10^{-3}$ ($-0.61$ / $-0.42\,\sigma_\text{reco}$, Fig.~\ref{fig:alphas}); the LAT value equals the earlier
gen--reco difference measured with unit TMD priors (z-data-unfolding study), so it does not come from those priors.
$\sigma(\alphas)$ is 3--8\,\% larger at gen. The gen fits want more CS damping at small $b$ ($\lambda_{2\nu}$ 0.079 vs 0.040 in
LAT) and less TMD $\lambda_2$ \taskref{261017-gen-cold-fits}.

\textbf{Physics read.} The nominal fits no longer depend on a warm-start history. With the NP $\lambda$ prior-free, the lattice is the only NP information and is
in $\approx3\sigma$ tension with the Z data at both levels; the cost of including it is $+5\,\%$ (reco) / $+1\,\%$ (gen) in $\sigma(\alphas)$
and a shift of about a quarter to a third of $\sigma$.

\subsection{The NP functions: the lattice moves damping from the CS kernel to the TMD factor}

\begin{figure}[htb]
  \centering
  \includegraphics[width=0.64\linewidth]{261015-cold-tc-fits/np_functions_NOLAT_vs_LAT/np_functions.pdf}\\[2pt]
  \includegraphics[width=0.64\linewidth]{261017-gen-cold-fits/np_gen_vs_reco_LAT/np_functions.pdf}
  \caption{Post-fit NP functions with 68\,\% bands from Gaussian toys of each fit's face-held covariance; dashed:
  prefit anchor. Left: CS rapidity anomalous dimension $\tilde\gamma_\nu^\text{NP}(b_T)$; right: TMD factor
  $F_\text{eff}(b_T,y)$ at $y=0$ and 2.5. Top: reco NOLAT vs reco LAT; bottom: gen LAT vs reco LAT. The grey/black
  ``lattice fit'' curve is the plotter's built-in reference (arXiv:2506.13874), \emph{not} the ASWZ $\chi^2$ used in LAT; NOLAT's $|Y|=2.5$ band runs away because toys around a point on the curved cubic face
  cross it: it is not a physical band. Bands pinch where a parameter sits on a face. \taskref{261015-cold-tc-fits} \taskref{261017-gen-cold-fits}}
  \label{fig:np}
\end{figure}

Without the lattice the CS kernel turns on late ($\lambda_{2\nu}=0$) and saturates early, while the TMD factor is weakly
damped at large $b_T$; with the lattice the CS kernel follows the lattice and the TMD factor damps much more strongly
(Fig.~\ref{fig:np}, top). The lattice thus decides which NP group dominates $\sigma(\alphas)$. Gen and reco LAT agree within the bands
(bottom); the gen CS kernel is slightly steeper at small $b_T$ and its TMD factor correspondingly less damped.

\subsection{\texorpdfstring{The \ptll\ misfit: a characterised mis-modelling, not adopted as a systematic}{The ptll misfit: a characterised mis-modelling}}

\caveat{All tests below are rabbit fits on the reco LAT configuration (NOLAT where stated), warm from LAT's minimum,
$\alphas$ shifts in units of LAT's post-fit $\sigma(\alphas)$. A first, surrogate-based pass of this task (quadratic
model of the saturated sub-fit, not rabbit) read the misfit as a low-order shape; that reading was wrong and is
superseded by the exact decomposition below \taskref{261017-ptll-gof-forensics}.}

\begin{figure}[htb]
  \centering
  \begin{subfigure}[t]{0.29\linewidth}\centering
    \includegraphics[width=\linewidth]{261015-cold-tc-fits/LAT/ptll/postfit_ptll_ch0_Project_ch0_ptll_preliminary.pdf}
    \caption{reco LAT post-fit \ptll}\end{subfigure}\hspace{0.03\linewidth}
  \begin{subfigure}[t]{0.35\linewidth}\centering
    \includegraphics[width=\linewidth]{261017-ptll-gof-forensics/p3/leading_mode_and_smooth_fit.pdf}
    \caption{the residual wave (LAT top, NOLAT bottom)}\end{subfigure}
  \caption{(a) Post-fit \ptll\ projection of the reco LAT fit; the label's ``$p=0\,\%$'' is 0.02\,\% rounded. (b) Post-fit
  residual (data$-$postfit)/postfit in \% with errors from the exact profiled residual covariance $R$; red: its component
  along the leading eigenmode of $R$ (a covariance-built shape, only its amplitude is from data); blue: best smooth
  7- and 10-parameter Legendre shapes. LAT's $R$ is approximate (it misses the lattice term); NOLAT's is validated to
  1\,\%. \taskref{261015-cold-tc-fits} \taskref{261017-ptll-gof-forensics}}
  \label{fig:ptll}
\end{figure}

The \ptll\ projection fails in every physical-NP fit since April 2026, across models, PDFs and cards; in these fits
$q=79.1/39$ (LAT) and 68.4/39 (NOLAT), and the gen $p_\mathrm{T}^V$ projection fails as well (Table~\ref{tab:fits}), which points to the
model side (the unfolded spectrum inherits the reco response, so a detector cause is not excluded; no detector template
was tested). The residuals are statistically ordinary bin by bin (data-only
$\sum z^2=41.5/39$ LAT, 39.7/39 NOLAT) and the misfit is common to all rapidities: splitting \ptll\ into 6 \yll\ bands adds
$\chi^2=210.6/195$ (LAT, $p$ 21\,\%) and 189.1/195 (NOLAT) \taskref{261017-ptll-gof-forensics}. It has two parts.

\textbf{(i) A smooth wave.} With the exact post-fit covariance the misfit is a $y$-independent $\pm0.2$--$0.3\,\%$ wave
between 3 and 22\,GeV ($-0.2\,\%$ at 5--6, $+0.3\,\%$ at 9--12, $-0.2\,\%$ at 16--20\,GeV; Fig.~\ref{fig:ptll}b), with
$<2\,\%$ of $\chi^2$ in flat to cubic components; one covariance mode with six sign changes carries 19--25 units, and after
10 smooth parameters the remainder is 35/29 (NOLAT, $p$ 20\,\%). The nested coarse-rebin saturated fits confirm it, and
put LAT's $\approx10$ extra units over NOLAT in this finer structure (remainder after 10 coarse bins 45.5/29 LAT, 34.7/29
NOLAT).

\textbf{(ii) A normalisation--shape tension.} The coarse saturated fits release mainly one overall normalisation
($\approx22$ units, equal coarse scales $\approx0.89$). Freeing the normalisation directly (luminosity unconstrained) is the
most informative test (Table~\ref{tab:ptlltests}): the data prefer $-12.3\pm3.0\,\%$, $2\Delta\mathrm{NLL}=16.1$ on 1 dof
($4.0\sigma$, asymptotic; Wilks is not guaranteed along this nonlinear $\alphas$--$\lambda_4$--normalisation valley), and $\alphas$ moves by $+4.2\,\sigma$ with $\sigma(\alphas)$ $\times1.53$ ($\lambda_4$ 0.072$\to$0.181). The
\ptll\ shape thus prefers an $\alphas$/NP point whose larger cross-section the luminosity (1.2\,\%), EW and FSR constraints
forbid; the nominal sits at a compromise held by those priors. An earlier, template-NP version of this test found
$6.1\sigma$. At gen level the post-fit $p_\mathrm{T}^V$ spectrum sits 3--4\,\% low, a coherent offset the correlated
unfolded covariance allows (linear $\chi^2$ $p$ 21\,\%) \taskref{261017-gen-cold-fits}. Even with the normalisation free the \ptll\ projection stays poor (62.96/38, $p$ 0.66\,\%): that is the wave.

\begin{table}[htb]
  \centering\small
  \caption{Decisive \ptll\ tests on reco LAT (reference: full 752.95/774, \ptll\ 79.11/39). ``Theory'' = A0 scale, FO
  envelope, $b/c$ masses, EW, FSR, luminosity priors. With the normalisation free the \ptll\ ndf is 38 (rabbit prints 39).
  \taskref{261017-ptll-gof-forensics}}
  \label{tab:ptlltests}
  \begin{tabular}{lcccc}
    \toprule
    test & $\Delta q_\text{full}$ & \ptll\ $q$/ndf, $p$ & $\dalphas$ [$\sigma_\text{LAT}$] & $\sigma(\alphas)$ ratio \\
    \midrule
    A0 scale prior $\times2$ & $-12.7$ & 67.54/39, 0.31\,\% & $-0.44$ & 1.00 \\
    all theory priors $\times2$ & $-22.5$ & 58.58/39, 2.3\,\% & $-0.01$ & 1.26 \\
    free normalisation & $-16.1$ & 62.96/38, 0.66\,\% & $+4.2$ & 1.53 \\
    \bottomrule
  \end{tabular}
\end{table}

Inflating the theory priors absorbs only part of the wave, and the pulls run with the prior: at doubled width the A0
15--20\,GeV scale and the FO-envelope nuisances still sit at $-1.9$ and $-2.1\sigma$ of the \emph{new} width ($-3.8$,
$-4.1\sigma$ in original units): a missing shape, not a too-narrow envelope;
$\alphas$ does not move ($-0.01\,\sigma$).

\textbf{Physics read and decision.} $\alphas$ is insensitive to the wave but depends, at the $4\sigma$ level, on the
absolute normalisation constraint. On 2026-10-20 Luca decided not to adopt any remedy (no normalisation-free systematic,
no shape nuisance): the \ptll\ misfit is reported as a known, characterised mis-modelling of the model's $q_T$ shape.
Untested: A0 modelling, $b/c$ masses in the resummation, QED, resummation order.

\subsection{The new fits are the same minima as the old trust-krylov + wall fits}

The earlier conclusions rest on warm trust-krylov fits with a penalty wall and unit priors on NP coordinates.
Evaluating the new objectives at the old points,
one Newton step from the old minima with the priors removed, and warm rabbit refits from the old points all land on the
new minima: the old points lie only 0.0023 above them, Newton predicts them to $\le8\times10^{-3}\,\sigma$, the refits
reach them to $\le5\times10^{-6}\,\sigma$, with the same faces. $\alphas$ moves by $\mp0.017\,\sigma$ (the prior removal), $\sigma(\alphas)$ by
$<0.2\,\%$, and the Z--lattice tension is 13.51 vs 13.48 before. The cold NOLAT found the lower of the two known Z-only
basins (the other sits 0.57 higher) \taskref{261017-minima-vs-trustkrylov}. The old conclusions stand.

% ============================================================================================
\section{Robustness, systematics and caveats}

\textbf{Saturated sub-fits.} Restarting scipy's barrier at $\mu_0=0.1$ from a minimum made the plain LAT \ptll\ sub-fit
crawl for 12.8\,h and stop on \code{xtol}; a converged rerun changed $q$ by $1.7\times10^{-4}$, so the quoted $q$ stand
\taskref{261016-saturated-subfit-diagnosis}. In Hessian-vector products the preconditioner alone gives $3.1\times$
(4667$\to$1504) and with the warm start $4.75\times$ (982), at the same point; an interim claim (2026-10-17) that most of
the gain came from the warm start was wrong: most of the speed comes from the preconditioner ($3.1\times$); the warm
start adds $1.5\times$ and fixes the false stops at releasing faces
\taskref{261015-trustconstr-precondition} \taskref{261015-trustconstr-warmstart}. The gen $p_\mathrm{T}^V$ sub-fits end on
\code{xtol} at optimality $\le2\times10^{-6}$ after $\approx5$\,h; their saturated covariances fail Cholesky ($q$ unaffected)
\taskref{261017-gen-cold-fits}. The ptll tests of Table~\ref{tab:ptlltests} ran at \code{gtol} $10^{-5}$, shown to move $q$ by $\lesssim10^{-6}$.

\textbf{Impacts.} Face-held traditional impacts match exact references and are bit-identical without faces \taskref{261016-face-held-impacts}. LAT's
impacts are those \emph{on} its weakly binding face: released, its total would be $1.28$ instead of $1.14\times10^{-3}$, so
the interval along that direction is one-sided.

% ============================================================================================
\section{What changed}

\begin{itemize}
  \item \textbf{rabbit (WMass):} \#194 spectral preconditioner under trust-constr (chain rule on the constraints; warm
  fits only) \emph{merged}; \#195 \code{--minimizerWarmStart} (\code{auto}/\code{on}/\code{off}; consistent slacks, centred $\mu_0$, a Newton
  release step at releasing faces) and \#196 face-held traditional impacts \emph{ready}, rebased on main, merge cleanly
  together \taskref{261020-pr-shepherd}.
  \item \textbf{WRemnants:} the CS $\lambda$ are prior-free by default (\code{cc3859f4}) and the two-panel NP plotter
  \code{scetlib_ad/np_function_plots.py} was ported (\code{74c7e3e3}), both on the closed \#715, which was split into
  reviewable PRs \taskref{261017-wremnants-pr-split}: \#736 (setupRabbit fixes) \emph{merged}; \#737 (response inputs;
  \code{--responseGenBinning} now defaults to \code{none}), \#738 (\code{ptZgen_binning_corr} edges, needs the
  theory-correction producers' sign-off) and \#739 (core param model, after \#737) \emph{ready}. The second wave (NP
  constraints, lattice term, validation tools, cache-to-theory-correction, NP plots, a rename of the response aux group to
  \code{response_matrix} with an old-name fallback) is built locally, not pushed \taskref{261020-pr-shepherd}.
  \item \textbf{WRemnantsHelpers:} new \code{--constraints} and \code{--lattice} options of \code{workflows/fitterAD.sh},
  and the NP-plot wrapper \code{workflows/plotNPFunctions.sh} (\code{010d37b}, \code{5c80345}, pushed); two validation scripts moved to
  \code{scripts/scetlib_ad/} (\code{d5508b5}, local). \textbf{knowledge/:} the AD post-fit recipe in
  \code{20_frameworks/nominal_workflow.md}; a trust-constr warm-start section in
  \code{20_frameworks/rabbit_minimizer_tolerances.md}.
\end{itemize}

% ============================================================================================
\section{Conclusions and open items}\label{sec:open}

The nominal CT18Z fits are cold, reproducible and equal to the earlier wall fits up to the intended prior removal.
Per configuration reco and gen share the NP faces; the $\approx3\sigma$ Z--lattice tension moves NP damping from the CS
kernel to the TMD factor and $\alphas$ up by 0.2--0.35$\,\sigma$; gen lies 0.4--0.6$\,\sigma$ below reco. The full fits are acceptable, but the model's $q_T$ shape is not: a $y$-independent few-per-mille
wave, which $\alphas$ ignores, plus a $4\sigma$ normalisation--shape tension to which $\alphas$ is sensitive; by decision
this is quoted as a characterised mis-modelling.

Open, none assigned: \textbf{(i)} both LAT fits sit on a weakly binding face, so an $\alphas$ interval for them needs a
likelihood scan; \textbf{(ii)} the NOLAT free-normalisation refit was not run;
\textbf{(iii)} merging \#195/\#196 and \#737--\#739, then opening the second WRemnants wave.

\end{document}
